\chapter{최소 반례는 3-연결이다}
\label{chap:reduce}

\section{붙이기}

세 보조정리 모두 회전 체계 수술이고, 면 수는 보조정리~\ref{lem:swap}로 센다. 2-합을 직접 증명하는
이유는 결정~D4(\ref{chap:plan}장)에 있다. 수는 서론의 세는 규약대로 고정된 \(V\)(\(n = |V|\)) 위에서
센다. \(F = O + I\)와 \(c\)에 고립점이 들어가므로 붙인 그래프의 \(F\), \(c\)는 그냥 더해지지 않는다.
세 경우 모두 \(I = I_1 + I_2 - n + |S_{G_1} \cap S_{G_2}|\)이다.

\begin{lemma}[서로소 합]
\label{lem:sum0}
\uses{def:IsPlanar}
\lean{Kuratowski.IsPlanar.sup_of_disjoint, Kuratowski.IsPlanar.sum,
  Kuratowski.RotationSystem.sum, Kuratowski.RotationSystem.sum_eulerChar}
\leanok
유한한 꼭짓점 타입 \(V\) 위의 \(G_1\), \(G_2\)가 평면이고 지지 집합이 서로소(\(S_{G_1} \cap S_{G_2} = \emptyset\))이면 \(G_1 \sqcup G_2\)는
평면이다. 이때 변 집합도 서로소다.
\end{lemma}

\begin{proof}
\uses{lem:map_iff, thm:mono}
\leanok
먼저 서로소 타입 합 \(V_1 \oplus V_2\) 위의 그래프에서 두 회전 체계를 나란히 놓는다. 다트는 두 다트
집합의 서로소 합이고, 회전과 면 순열은 각 쪽에서 독립적으로 작용한다. 따라서 면 순환·변·고립점·연결
성분 수가 모두 더해지며 \(\chi = \chi_1 + \chi_2 = 2(c_1 + c_2)\)다. 이것이 \texttt{IsPlanar.sum}이다.

원래 고정 타입 \(V\)의 두 그래프를 각각 지지 집합으로 제한한다. 고립점은 보조정리~\ref{lem:map_iff}로
제거할 수 있다. 두 지지 집합이 서로소이므로 그 타입 합에서 \(V\)로 가는 포함 사상이 단사다. 위 평면
그래프를 이 단사 사상으로 옮기면 \(G_1 \sqcup G_2\)를 포함하므로 평면성의 단조성으로 끝난다.
예: \(V = \mathrm{Fin}\,4\), \(G_1 = \{01\}\), \(G_2 = \{23\}\)이면 고립점을 포함한 고정 타입의 수는
\(F = 3 + 3 - 4 = 2\), \(c = 3 + 3 - 4 = 2\)다.
\end{proof}

\begin{lemma}[한 점에서 붙이기]
\label{lem:sum1}
\uses{def:IsPlanar}
\lean{Kuratowski.IsPlanar.sup_of_inter_eq_singleton, Kuratowski.RotationSystem.glue,
  Kuratowski.RotationSystem.orbitCount_glue_facePerm,
  Kuratowski.card_connectedComponent_sup_of_inter_eq_singleton,
  Kuratowski.ncard_isolated_sup_of_inter_eq_singleton}
\leanok
유한한 꼭짓점 타입 \(V\) 위의 \(G_1\), \(G_2\)가 평면이고 지지 집합이 꼭짓점 하나에서만 만나면(\(S_{G_1} \cap S_{G_2} = \{v\}\))
\(G_1 \sqcup G_2\)는 평면이다. 이때 공유하는 변은 없다 (공유하는 변의 두 끝점은 서로 다르고 둘 다
\(S_{G_1} \cap S_{G_2}\)에 있어야 한다).
\end{lemma}

\begin{proof}
\uses{lem:swap, lem:map_iff}
\leanok
\(v\)에서 두 회전 순환을 이어 붙인다. \(G_1\)의 \(v\) 다트 \(p\)와 \(G_2\)의 \(v\) 다트 \(q\)에 대해 회전에
호환 \((p\ q)\)를 곱하면 서로 다른 두 순환이 하나가 된다. 면 순열에는 호환 \((\bar p\ \bar q)\)가 곱해지고
(\(\bar p\)는 \(p\)의 반대 다트), \(\bar p\)와 \(\bar q\)는 서로 다른 그래프의 면에 있으므로 두 면이 하나로
합쳐진다. 따라서 \(O = O_1 + O_2 - 1\)이다. 또 \(E = E_1 + E_2\), \(I = I_1 + I_2 - n + 1\)이므로
\(F = F_1 + F_2 - n\)이고, \(v\)를 지나는 두 성분이 하나가 되어 \(c = c_1 + c_2 - n\)이다.
이 성분 수 등식은 합집합의 성분을 (i) \(G_1\)의 지지 집합에서의 성분, (ii) \(G_2\)의 지지 집합에서
\(v\)의 성분을 뺀 나머지 성분, (iii) 두 그래프의 공통 고립점으로 분류하는 일대일 대응으로 증명한다.
고정 타입과 지지 집합 사이의 성분 수는 보조정리~\ref{lem:map_iff}에서 얻는다. Lean에서는 자연수 뺄셈 대신
\(O + 1 = O_1 + O_2\), \(I + n = I_1 + I_2 + 1\), \(c + n = c_1 + c_2\)를 쓴다. 그러므로
\(\chi = \chi_1 + \chi_2 - 2n = 2c\)다.
\end{proof}

\begin{lemma}[공통 성분에 놓인 교집합의 세기]
\label{lem:shared_component_counts}
\uses{lem:map_iff}
\lean{Kuratowski.card_connectedComponent_sup_of_shared_component,
  Kuratowski.card_connectedComponent_sup_of_inter_eq_edge, Kuratowski.ncard_isolated_sup_add}
\leanok
유한한 \(V\) 위의 \(G_1,G_2\)에서 \(v\in S_{G_1}\cap S_{G_2}\)이고, 모든 공통 지지점이 각 그래프의
지지 집합에 제한한 그래프에서 \(v\)와 같은 성분에 있다고 하자. 고립점을 포함한 성분 수에 대해
\[
 c(G_1\sqcup G_2)+n+1=c(G_1)+c(G_2)+|S_{G_1}\cap S_{G_2}|
\]
이다. 특히 지지 집합 교집합이 \(\{x,y\}\)이고 두 그래프 모두에 변 \(xy\)가 있으면
\(c(G_1\sqcup G_2)+n=c(G_1)+c(G_2)+1\)이다. 고립점 수에는 추가 연결 가정 없이 항상
\(I(G_1\sqcup G_2)+n=I(G_1)+I(G_2)+|S_{G_1}\cap S_{G_2}|\)가 성립한다.
\end{lemma}

\begin{proof}
\leanok
합집합의 성분을 \(G_1\) 지지 집합의 성분, \(G_2\) 지지 집합에서 \(v\)의 성분을 뺀 나머지 성분,
두 그래프의 공통 고립점으로 분류한다. 공통 지지점들이 양쪽에서 같은 지정 성분에 있으므로 이 분류는
변을 따라 변하지 않고 성분과 일대일 대응한다. 보조정리~\ref{lem:map_iff}의 지지 집합 제한 공식과
포함 배제로 식을 얻는다. 고립점 식은 지지 집합 합집합의 여집합을 세면 된다.
\end{proof}

\begin{lemma}[순열에서 두 점 삭제의 순환 수]
\label{lem:pair_deletion}
\uses{lem:swap, lem:orbitCount_cut}
\lean{Kuratowski.PermutationDeletion.orbitCount_cut_pair_face}
\leanok
유한한 타입 \(D\)(같음 판정 가능)의 순열 \(\rho,\sigma\)와 서로 다른 \(a,b\in D\)가
\(\sigma(a)=b\), \(\sigma(b)=a\)를 만족한다고 하자. 또
\[
 \rho(a)\ne a,\qquad \rho(b)\ne b,\qquad \rho(b)\ne a
\]
이고 \(a,b\)가 \(\rho\sigma\)의 서로 다른 순환에 있다고 하자.
\(\psi=\rho\sigma(a\ b)\)에서 \(a\), \(b\)를 차례로 잘라 고정점으로 만든 뒤 두 점을 뺀 부분형에
제한한 순열을 \(\psi'\)라 하면, \(O(\psi')+1=O(\rho\sigma)\)다.
이 정리는 단순 그래프의 다트라는 가정을 쓰지 않는다.
\end{lemma}

\begin{proof}
\leanok
\(\psi(a)=\rho(a)\), \(\psi(b)=\rho(b)\)다. 위 세 부등식 때문에 두 잘라 내기가 각각 고정점 순환을
하나씩 추가한다. 부분형에 제한하면 이 두 고정점 순환이 제거되므로 순환 수는 \(O(\psi)\)와 같다.
한편 서로 다른 두 순환에 호환 \((a\ b)\)를 곱하면 두 순환이 합쳐지므로
\(O(\psi)+1=O(\rho\sigma)\)다.
\end{proof}

\begin{lemma}[변 하나를 따라 붙이기]
\label{lem:sum2}
\uses{def:IsPlanar}
\lean{Kuratowski.IsPlanar.sup_of_inter_eq_edge, Kuratowski.GlueEdge.rotation,
  Kuratowski.GlueEdge.dartEquiv, Kuratowski.GlueEdge.joinedFace_bigon,
  Kuratowski.GlueEdge.rotation_facePerm, Kuratowski.GlueEdge.orbitCount_joinedFace,
  Kuratowski.GlueEdge.orbitCount_rotation_facePerm, Kuratowski.GlueEdge.ncard_edgeSet_sup}
\leanok
유한한 꼭짓점 타입 \(V\) 위의 \(G_1\), \(G_2\)가 평면이고 지지 집합이 두 점에서만 만나며
(\(S_{G_1} \cap S_{G_2} = \{x, y\}\), \(x \ne y\)) \(xy\)가 두 그래프 모두의 변이면
\(G_1 \sqcup G_2\)는 평면이다. 이때 공유하는 변은 \(xy\) 하나뿐이다.
\bfcheck{\texttt{[sum2-수술]}: 공통 변을 \(01\)로 고정한 꼭짓점 5개 이하의 모든 해당 그래프 쌍과,
공통 고립점이 있는 두 \(K_4\) 예를 합쳐 서로 다른 \((n,E_1,E_2)\) 기준 그래프 쌍 1,617개·회전 쌍 49,664개.
같은 변 집합을 만드는 꼭짓점 배정은 중복을 뺀다. 원래 회전이 종수 0이
아닌 경우도 아래 두각형과 면 수 식을 검사한다. 원래 회전이 둘 다 종수 0이면 결과 회전도 종수 0이다.
\texttt{[sum2-수술-음성]}: \(x\), \(y\) 쪽 선행 다트 선택을 바꾸면 각각 25,064개에서 두각형 또는 면 수
조건이 실패한다. 유한 검사이며 일반 증명을 대신하지 않는다.}
\end{lemma}

\begin{proof}
\uses{lem:swap, lem:orbitCount_cut, lem:pair_deletion, lem:shared_component_counts, lem:sum0}
\leanok
두 그래프의 다트에 출처를 붙인 서로소 합 \(D = D_1 \oplus D_2\)에서 작업한다. 공통 변의 두 사본을
일단 유지하므로 이 중간 상태는 단순 그래프의 회전 체계가 아니라 \(D\) 위의 순열이다. 역다트 사상을
\(\sigma\), 원래 회전의 서로소 합을 \(\rho\), 면 순열을 \(\phi = \rho\sigma\)라 하자.
\(p \in D_1\), \(q \in D_2\)는 각각 \(x \to y\)인 다트다. 다음과 같이 둔다.
\[
 a = \rho_2^{-1}(q), \qquad b = \rho_1^{-1}(\bar p), \qquad
 \rho_0 = \rho\,(p\ a)\,(b\ \bar q).
\]
\(p,a\)는 \(x\)에서, \(b,\bar q\)는 \(y\)에서 출발하므로 두 호환은 서로소다. 각 꼭짓점의 두 회전
순환을 합치며, 이 선행 다트 선택은 다음 두각형을 보장한다.
\[
 \phi_0(\bar p) = \rho_0(p) = q, \qquad
 \phi_0(q) = \rho_0(\bar q) = \bar p.
\]
\begin{itemize}
  \item 면 순열로 옮기면 \(\phi_0 = \phi\,(\bar p\ \bar a)\,(\bar b\ q)\)다. 첫 호환은 서로 다른
        그래프의 면을 합친다. \(\phi(\bar b)=\bar p\), \(\phi(\bar a)=q\)이므로 둘째 호환의 두 다트는
        합쳐진 그 면에 있다. 둘째 호환은 이를 나눈다. 따라서 \(O_0 = O_1 + O_2\)다.
  \item \(\rho_0\)에서 \(q,\bar q\)를 차례로 잘라 내고 두 점을 뺀 부분형에 제한한다. 남은 다트는
        \(G_1 \sqcup G_2\)의 다트와 일대일로 대응한다. 두 호환이 \(x\), \(y\)에서 각각 회전 순환을
        합쳤고, 잘라 내기는 남은 다트 사이의 같은 순환 관계를 보존한다. 따라서 부분형으로 제한하고
        다트 대응으로 옮긴 뒤에도 각 꼭짓점의 다트는 한 순환을 이룬다.
  \item \(q\)의 면은 정확히 \(\{\bar p,q\}\)이고 \(\bar q\)는 다른 면에 있다. 삭제의 면 순열 식
        \(\mathrm{cut}(\mathrm{cut}(\phi_0(q\ \bar q),q),\bar q)\)를 일반 순열에서 증명한다.
        두 면을 합친 뒤 생긴 두 고정점을 부분형에서 빼면 \(O+1=O_0=O_1+O_2\)다.
\end{itemize}
차수가 1인 경우에는 \(a=q\) 또는 \(b=\bar p\)일 수 있지만 두 호환의 출발 꼭짓점이 다르므로 위 계산은
그대로 성립한다. 거울상을 선택하는 대신 선행 다트를 명시해 방향을 고정했다.

세는 규약으로는 \(E+1=E_1+E_2\), \(I+n=I_1+I_2+2\), \(c+n=c_1+c_2+1\)이다. 공통 변 때문에 두
공통 꼭짓점은 각 그래프에서 같은 성분에 속하므로 지지 집합의 성분 하나씩만 합쳐진다. 따라서
\(F=F_1+F_2-n+1\), \(\chi=\chi_1+\chi_2-2n+2=2c\)다. 두 삼각형을 붙인 마름모는 \(O=3,F=3\)이다.

Lean에서는 \texttt{GlueEdge.lean}이 이중 다트 위의 순열에서 실제 합집합의 회전 체계까지 구성한다.
\texttt{PermutationDeletion.lean}의 일반 순열판을 중간 상태에 적용하고,
\texttt{GlueOne.lean}의 고립점·성분 수 공식을 써서 마지막 오일러 식을 얻는다.
최종 평면성 정리는 유한한 \(V\)만 가정하며, 별도의 결정 가능성 인스턴스를 요구하지 않는다.
\end{proof}

\section{3-연결성과 수축}

\begin{definition}[\(k\)-연결]
\label{def:conn3}
\lean{Kuratowski.ConnectedOn, Kuratowski.TwoConnected, Kuratowski.ThreeConnected,
  Kuratowski.deleteVertices, Kuratowski.connectedOn_of_clique, Kuratowski.not_connectedOn_of_isolated}
\leanok
유한한 꼭짓점 타입 \(V\) 위에서 꼭짓점 집합 \(A\)가 \(G\)에서 \emph{\(k\)-연결}이라는 것은
\(|A| \ge k + 1\)이고, \(|S| < k\)인 꼭짓점 집합 \(S\)를 지워도
\(A \setminus S\)의 점들이 \(G - S\)에서 서로 이어져 있다는 것이다. 양의 \(k\)에서는
이것이 \(k - 1\)개 이하 삭제와 같다. \(k = 0\)에서는 \(A\)가 비지 않는 크기 조건만 남는다. 여기서
\(G - S\)는 \(S\)에 닿는 변을 모두 지운 그래프다 (꼭짓점 타입은 그대로 \(V\)). \(G\)가 \(k\)-연결이라는 것은
지지 집합 \(S_G\)가 \(G\)에서 \(k\)-연결이라는 것이다. 2-연결과 3-연결은 \(k = 2, 3\)인 경우다
(\texttt{TwoConnected}, \texttt{ThreeConnected}).

\(A \setminus S\)는 지우기 \emph{전}의 집합에서 \(S\)만 뺀 것이다. 지운 뒤 \(G - S\)의 지지 집합을 다시
계산하면 안 된다. 그렇게 하면 별 \(K_{1,3}\)의 중심을 지웠을 때 지지 집합이 비어 공허하게 이어진 것이
되고, 별이 2-연결이 되어 버린다. 모든 그래프가 고정된 꼭짓점 타입 \(V\) 위에 있고 수축이 고립점을
만들기 때문에 지지 집합을 기준으로 정의한다 (결정~D2).
\bfcheck{\texttt{[conn]}: 별 \(K_{1,3}\), 삼각형과 고립점, \(K_4\)와 고립점, 길이 2인 경로,
보조정리~\ref{lem:J2}의 반례 그래프에서 이 정의대로 판정한다. 다른 검사도 모두 이 정의를 쓴다.}
같은 경계 예제는 \texttt{ConnectedExamples.lean}에서도 Lean 커널로 확인한다. 클리크의 연결성이나
삭제 뒤 고립된 점을 직접 사용하므로, Python 판정 결과를 가정으로 삼지 않는다.
\end{definition}

\begin{definition}[변 수축]
\label{def:contract}
\lean{Kuratowski.contractEdge, Kuratowski.contractEdge_adj, Kuratowski.isIsolated_contractEdge,
  Kuratowski.contractEdge_adj_left, Kuratowski.neighborSet_contractEdge_left,
  Kuratowski.contractEdge_adj_of_ne, Kuratowski.contractEdge_self_adj,
  Kuratowski.contractEdge_top_zero_one}
\leanok
\(G/xy\)는 \(y\)를 \(x\)에 합친 그래프다 (\(x \sim y\)). \(p \ne q\)이고 \(p, q \ne y\)일 때 \(p \sim q\)는
원래 인접이거나, \(p\)(또는 \(q\))가 \(x\)이고 다른 쪽이 \(y\)의 이웃인 경우다. \(p \ne q\) 조건이 없으면
\(x\)가 \(y\)의 이웃이라서 고리 \(xx\)가 생긴다. \(y\)는 고립점이다. 겹친 변은 하나로 친다 (단순 그래프).
\(x \ne y\)이면 \(x\)의 새 이웃은 \((N(x) \cup N(y)) \setminus \{x, y\}\)이고, \(x\), \(y\)가 아닌 두 점
사이의 인접은 그대로다. 예를 들어 \(K_4\)에서 변 \(01\)을 수축하면 \(1\)은 고립점이고 \(0\), \(2\), \(3\)은
삼각형을 이룬다.

Lean 정의는 모든 \(x\), \(y\)에 대해 주어진다. \(x = y\)이면 \(x\)에 닿는 변을 모두 지운 그래프가 되는데,
실제 수축은 \(x \sim y\)(따라서 \(x \ne y\))를 가정하므로 이 경우는 쓰이지 않는다.
\texttt{scripts/check\_blueprint.py}의 \texttt{contract}와는 \(x \ne y\)이면 늘 같고, \(x = y\)이면 일반적으로
다르다 (Lean은 \(x\)에 닿는 변을 지우고 Python은 \(G\)를 그대로 둔다). 변 수, 고립점 수, 성분 수는
아래 보조정리~\ref{lem:contract_counts}에서 조건을 구분해 다룬다.
\end{definition}

\begin{lemma}[수축의 변과 성분 수]
\label{lem:contract_counts}
\uses{def:contract}
\lean{Kuratowski.ncard_edgeSet_contractEdge_lt, Kuratowski.card_connectedComponent_contractEdge,
  Kuratowski.ncard_isolated_contractEdge, Kuratowski.ncard_edgeSet_of_dartEquiv_delete_pair,
  Kuratowski.le_contractEdge_of_isIsolated, Kuratowski.deleteVertices_contractEdge_left}
\leanok
유한한 꼭짓점 타입 위에서 \(x \sim y\)이면 \(E(G/xy) < E(G)\), \(c(G/xy) = c(G)+1\)이다.
수축 뒤 \(x\)가 고립점이 아니면 \(I(G/xy) = I(G)+1\)이다. \(xy\)가 \(K_2\) 성분이면 수축 뒤
\(x\), \(y\)가 모두 고립점이 되므로 마지막 식의 비고립 가정은 필요하다.
또 임의의 그래프 \(H\)와 다트 \(p \in D_G\)에 대해
\(D_G \setminus \{p,\bar p\} \simeq D_H\)인 전단사가 있으면 \(E(G)=E(H)+1\)이다.
\(H \le G\)이고 \(H\)에서 \(y\)가 고립점이면 \(H \le G/xy\)이고, 인접 가정 없이도
\((G/xy)-x = G-\{x,y\}\)다.
\end{lemma}

\begin{proof}
\leanok
\(y \mapsto x\)로 변을 보내면 수축 뒤의 변은 원래 변에서 \(xy\)를 뺀 집합의 상에 포함되므로 변 수는 줄어든다.
\(y\)가 아닌 두 점은 수축 전후에 이어짐이 같고, \(y\)는 수축 뒤 혼자 한 성분이다. 따라서 수축 뒤 성분은
원래 성분들과 \(y\)의 단독 성분으로 분류된다. \(x,y\) 밖 점들의 고립 여부는 유지된다.
두 다트를 뺀 전단사에서는 다트 수가 2 줄고 변 하나당 다트가 둘이므로 변 수 식이 따른다.
부분 그래프의 고립점에는 변이 닿지 않으므로 수축해도 그 부분 그래프의 변은 남는다.
마지막 그래프 등식은 인접 관계를 펼치면 바로 따른다.
\end{proof}

\begin{lemma}[연결도의 기본 성질과 최소 분리 집합]
\label{lem:connectivity_tools}
\uses{def:conn3}
\lean{Kuratowski.ConnectedOn.mono, Kuratowski.ConnectedOn.deleteVertex,
  Kuratowski.ConnectedOn.ncard_neighborSet_ge, Kuratowski.exists_adj_mem_sepComponent_of_connectedOn}
\leanok
유한한 꼭짓점 타입에서 \(A\)가 \(G\)에서 \(k\)-연결이면 \(j \le k\)인 \(j\)-연결이고,
각 \(x \in A\)의 이웃 수는 \(k\) 이상이다. \(A\)가 \((k+1)\)-연결이면 \(A \setminus \{v\}\)는
\(G-v\)에서 \(k\)-연결이다.
또 지지 집합이 \(k\)-연결이고 크기 \(k\)인 \(T\)가 지지 집합의 두 점을 분리하면, \(T\)의 각 점은
\(G-T\)의 각 지지 성분에 이웃을 갖는다. 정확히, \(c \in S_G \setminus T\), \(t \in T\)마다 어떤
\(d \in S_G \setminus T\)가 있어 \(c\)와 \(d\)가 \(G-T\)에서 이어지고 \(t \sim d\)다.
\end{lemma}

\begin{proof}
\leanok
삭제할 점 수의 하한을 낮추면 단조성이 따른다. 한 점 삭제 뒤에는 추가 삭제 집합에 그 점을 더해
원래 연결성을 쓴다. 이웃 수가 \(k\)보다 적으면 이웃을 모두 지워 \(x\)를 고립시키며, 크기 조건으로
기준 집합의 다른 점이 남아 모순이다. 마지막 주장에서 \(t\)가 한 성분에 이웃이 없으면 \(T \setminus \{t\}\)를
지워도 그 성분을 벗어나지 못한다. 크기가 \(k-1\)인 분리 집합이 되어 \(k\)-연결성에 모순이다.
\end{proof}

\begin{theorem}[수축 가능한 변]
\label{thm:contractible}
\uses{def:conn3, def:contract}
\lean{Kuratowski.ThreeConnected.exists_contractible,
  Kuratowski.exists_separates_of_not_threeConnected, Kuratowski.exists_adj_mem_sepComponent}
\leanok
유한한 꼭짓점 타입 위에서 지지 집합이 5개 이상인 3-연결 그래프에는 \(G/xy\)가 3-연결인 변 \(xy\)가 있다.
\bfcheck{\texttt{[T]}: 지지 집합이 5개인 이름표 붙은 3-연결 그래프 전부와 바퀴 \(W_5\), \(W_6\), 삼각기둥,
정팔면체, 정육면체 (모두 31개)}
\end{theorem}

\begin{proof}
\leanok
Diestel, \emph{Graph Theory}, 6판 보조정리 3.2.4(5판 3.2.1)의 증명을 지지 집합 기준으로 옮긴다.
원문은 \url{https://www.math.uni-hamburg.de/home/diestel/books/graph.theory/preview/Ch3.pdf}에 있다.
유한한 \(V\)에서 논증하며, 삭제 후에도 지우기 전 지지 집합에서 살아남은 점들을 대상으로 한다.
\begin{enumerate}
  \item 3-연결성과 \(|S_G|\ge5\)에서 \((G/xy).\mathrm{support}=S_G\setminus\{y\}\)이고 크기가 4 이상이다.
        수축 가능한 변이 없으면 각 변 \(xy\)에 대해 \(\{x,y,z\}\)가 분리 집합인 \(z\)를 얻는다.
        수축 그래프의 작은 분리 집합 \(S\)가 \(x\)를 포함하지 않으면 \(G\)의 걸음을 \(y\mapsto x\)로
        사영해 모순을 얻는다(합쳐지는 \(xy\) 변은 건너뛴다). \(x\)를 포함하면 \(S\cup\{y\}\)가
        \(G\)를 분리하며, 지지 집합과의 교집합이 정확히 \(\{x,y,z\}\)가 된다.
  \item 분리 집합 \(T=\{x,y,z\}\)의 진부분집합은 분리할 수 없다. 따라서 \(T\)의 각 점은
        \(G-T\)의 각 지지 성분에 이웃을 갖는다. 그 가운데 성분 \(C\)의 크기가 최소인 것을 고른다.
  \item \(z\)의 이웃 \(v\in C\)를 고르고 변 \(zv\)에 같은 논증을 적용해 분리 집합
        \(T'=\{z,v,w\}\)를 얻는다. \(x,y\) 중 \(T'\) 밖에 있는 점을 \(r\)로 고른다.
        둘 다 남으면 변 \(xy\)로 이어져 있다. \(r\)의 성분이 아닌 성분 \(D\)를 고르면 \(D\)는
        \(x,y\)를 모두 피한다.
  \item \(v\)는 \(D\) 안에 이웃 \(d\)를 갖는다. \(d\notin\{x,y,z\}\)이므로 \(d\in C\)다.
        \(D\) 안의 걸음은 \(\{x,y,z\}\)를 피하므로 \(D\subseteq C\)이고, \(v\in C\setminus D\)이므로
        \(|D|<|C|\)다. 이는 성분 크기의 최소성에 모순이다.
\end{enumerate}
\end{proof}

\begin{lemma}[들어 올리기]
\label{lem:lift}
\uses{def:contract, def:subdiv, def:K5K33}
\lean{Kuratowski.kuratowskiFree_contractEdge, Kuratowski.KuratowskiFree,
  Kuratowski.KuratowskiFree.contractEdge,
  Kuratowski.containsSubdivision_K5_or_K33_of_contractEdge,
  Kuratowski.containsSubdivision_K33_of_contractEdge,
  Kuratowski.SubdivisionPaths.containsSubdivision_K33_of_split}
\leanok
\(x \sim y\)이고 \(G\)가 \(K_5\)와 \(K_{3,3}\)의 세분을 포함하지 않으면 (KF) \(G/xy\)도 KF다.
Lean의 \texttt{KuratowskiFree G}는 이 두 세분이 없다는 조건의 논리곱이다. 이 보조정리에는 \(V\)의
유한성이나 별도의 결정 가능성 가정을 요구하지 않는다.
\end{lemma}

\begin{proof}
\uses{lem:subdivision_paths_forward, lem:paths_finite_source}
\leanok
대우를 보인다. 수축 그래프 \(G/xy\)에서 \(x\)를 남기고 \(y\)는 고립점으로 둔다. 그 안의 세분에서 경로
체계를 얻는다. \(K_5\), \(K_{3,3}\)에는 고립점이 없으므로 \(y\)는 분기점이 아니고 어느 경로에도 없다.
\begin{itemize}
  \item \(x\)가 분기점이 아니면 각 변을 원래 변 또는 \(y\)를 거치는 두 변으로 들어 올리고, 반복점이 있으면
        \texttt{bypass}로 경로로 줄인다. 새 내부는 원래 내부와 \(y\)에 포함된다. \(x\)를 지나던 무방향 변의
        경로는 많아야 하나이므로 \(y\)도 많아야 한 경로에만 추가된다. \(x\)가 쓰이지 않는 경우도 포함한다.
  \item \(x\)가 분기점이면 그 점에서 나가는 경로들의 첫 변을 본다. \(t \in \{x,y\}\)에 첫 변이 닿지 않는
        경로가 많아야 하나인 쪽을 새 분기점으로 고른다. 그 경로에만 다른 점을 끼우면 되므로 내부 서로소성이
        유지된다. 이는 4--0, 3--1 배치와 차수 3에서의 다수 쪽 선택을 포함한다.
  \item 어느 쪽을 골라도 나쁜 경로가 둘 이상이면, 첫 변이 \(x\)에만 닿는 두 경로 \(a_1,a_2\)와 \(y\)에만
        닿는 두 경로 \(b_1,b_2\)가 서로 다르다. 차수 3인 \(K_{3,3}\)에서는 불가능하다. \(K_5\)에서는 2--2
        분할이다. 분기점을 \(\{x,f(b_1),f(b_2)\}\), \(\{y,f(a_1),f(a_2)\}\)로 놓고 변 \(xy\), 두 경로
        \(x \to f(a_j)\), 두 경로 \(f(b_i) \to y\), 원래 경로 \(b_i\)--\(a_j\) 넷으로 \(K_{3,3}\)의
        경로 체계를 만든다. 원래 교차 경로들은 \(x\)를 피하므로 새 점이 추가되지 않는다.
\end{itemize}
마지막에는 \(K_5\) 또는 \(K_{3,3}\) 쪽 꼭짓점 타입이 유한한 경로 체계에서 세분을 구성한다. 따라서 \(K_{3,3}\)의 세분은 그대로
들어 올려지고, \(K_5\)의 세분은 \(K_5\)나 \(K_{3,3}\)의 세분으로 들어 올려진다. 둘 다 없다는 가정에 모순이다.
\end{proof}

\begin{lemma}[바깥 경로를 한 변으로 바꾸기]
\label{lem:path_replacement}
\uses{def:subdiv, def:K5K33, def:paths}
\lean{Kuratowski.KuratowskiFree.sup_edge_of_path, Kuratowski.KuratowskiFree.mono,
  Kuratowski.contractAvoidWalk, Kuratowski.contractAvoidWalk_length,
  Kuratowski.contractAvoidWalk_support}
\leanok
KF 그래프 \(G\), 부분 그래프 \(H \le G\), 서로 다른 \(x,z\)를 잇는 \(G\)의 단순 경로 \(P\)가 있고,
\(P\)의 내부 꼭짓점이 모두 \(H\)의 고립점이면 \(H+ xz\)도 KF다. 꼭짓점 타입의 유한성은 필요 없다.
\end{lemma}

\begin{proof}
\uses{lem:lift, lem:contract_counts}
\leanok
경로 길이에 대한 강한 귀납법. 한 변 경로면 \(H+xz \le G\)다. 첫 내부점 \(u\)가 있으면 \(xu\)를
수축한다. 보조정리~\ref{lem:lift}로 KF가 유지되고, \(u\)가 \(H\)의 고립점이라 \(H\)의 변도 남는다.
경로의 꼬리는 \(x\)를 피하므로 \(u \mapsto x\) 사영에서 접히는 변 없이 길이가 그대로다.
\texttt{bypass}로 경로로 줄이면 길이가 엄밀히 작아지고, 내부는 여전히 \(H\)의 지지 집합 밖이다.
귀납 가정으로 결론을 얻는다. KF의 부분 그래프 닫힘은 세분의 포함 관계에서 바로 따른다.
\end{proof}

\begin{theorem}[최소 반례는 3-연결]
\label{thm:reduce}
\uses{def:conn3, def:IsPlanar, def:subdiv, def:K5K33}
\lean{Kuratowski.isPlanar_or_threeConnected, Kuratowski.isPlanar_or_twoConnected,
  Kuratowski.separation_sup, Kuratowski.separation_support_inter_subset,
  Kuratowski.exists_path_in_sepComponent, Kuratowski.ncard_sup_edge_lt_of_two_missing}
\leanok
유한한 꼭짓점 타입 \(V\) 위에서 \(G\)가 KF이고, 같은 \(V\) 위의 변이 더 적은 모든 KF 그래프가 평면이면,
\(G\)는 평면이거나 3-연결이다.
\end{theorem}

\begin{proof}
\uses{lem:small, lem:sum0, lem:sum1, lem:sum2, lem:path_replacement, lem:connectivity_tools, thm:mono}
\leanok
\begin{enumerate}
  \item 지지 집합이 4개 이하이면 보조정리~\ref{lem:small}로 끝난다.
  \item 연결도가 부족하면 삭제 집합을 지지 집합과 교차시켜 \(T \subseteq S_G\)인 작은 분리 집합을 얻는다.
        \(G-T\)의 지지 성분 하나를 \(C\)라 하고, \(C\cup T\)를 남긴 그래프 \(G_1\)과 \(C\)를 지운
        그래프 \(G_2\)로 나눈다. 둘의 합집합은 \(G\)이고 지지 교집합은 \(T\)에 포함된다.
        \(|T| \le 1\)이면 각 쪽은 반대쪽 비고립점의 변을 하나 이상 놓친다. 따라서 변이 더 적고 부분 그래프라
        KF이므로 귀납 가정으로 평면이다. 보조정리~\ref{lem:sum0}, \ref{lem:sum1}로 합집합이 평면이다.
        이로써 비평면인 경우 2-연결이다.
  \item 3-연결이 아니면 \(|T|<3\)인 분리 집합이 있고, 2-연결성에서 \(|T|=2\)다. \(T=\{x,y\}\)라 하자.
        양쪽에 \(xy\)를 더한 \(G_1+xy\), \(G_2+xy\)를 본다.
    \begin{itemize}
      \item 보조정리~\ref{lem:connectivity_tools}에서 \(x\), \(y\)는 반대쪽 각 지지 성분에 이웃을 갖는다.
            그 성분 안의 걸음에 양 끝 변을 붙이고 \texttt{bypass}로 줄이면 \(x\)--\(y\) 경로를 얻는다.
            그 내부는 해당 쪽 그래프에서 고립점들이므로 보조정리~\ref{lem:path_replacement}로 추가 변 뒤에도 KF다.
      \item 각 쪽은 반대 성분으로 가는 서로 다른 두 원래 변을 놓치며 \(xy\) 하나만 더한다. 따라서 변 수는
            원래보다 엄밀히 적다. 두 변이 서로 다름은 반대 성분의 점이 \(\{x,y\}\) 밖에 있다는 데서 따른다.
    \end{itemize}
  \item 귀납 가정으로 두 추가 그래프는 평면이다. 지지 교집합은 정확히 \(\{x,y\}\)이고 \(xy\)가 양쪽에
        있으므로 보조정리~\ref{lem:sum2}를 적용한다. 원래 \(G\)는 그 합집합의 부분 그래프이므로
        정리~\ref{thm:mono}로 평면이다.
\end{enumerate}
\end{proof}
